% ECONOMICS_PROBLEM: {"id": "I15-474", "title": "Health–development feedbacks", "jel_code": "I15", "source_id": "OP-0037"}
% !TeX program = xelatex
\documentclass[11pt,a4paper]{article}
\usepackage[margin=23mm]{geometry}
\usepackage{fontspec}
\setmainfont{Latin Modern Roman}
\usepackage{amsmath,amssymb,mathtools}
\usepackage{xcolor,enumitem,booktabs,longtable,array,xurl,hyperref}
\definecolor{muted}{HTML}{53636C}
\hypersetup{colorlinks=true,urlcolor=blue,linkcolor=blue}
\setlength{\parindent}{0pt}
\setlength{\parskip}{5pt}
\newcommand{\R}{\mathbb R}
\newcommand{\E}{\mathbb E}
\newcommand{\N}{\mathbb N}
\newcommand{\ind}{\mathbf 1}
\newcommand{\Var}{\operatorname{Var}}
\newcommand{\Cov}{\operatorname{Cov}}
\newcommand{\supp}{\operatorname{supp}}
\DeclareMathOperator*{\argmin}{arg\,min}
\DeclareMathOperator*{\argmax}{arg\,max}
\newcommand{\entrymeta}[3]{{\sffamily\small\color{muted}\textbf{\texttt{#1}}\quad|\quad #2\par #3\par}}
\newcommand{\Problemref}[1]{\texttt{#1}}
\newcommand{\JELref}[2]{\texttt{#2} (JEL \texttt{#1})}
\begin{document}
\section*{Health–development feedbacks}
\textbf{Problem ID:} \texttt{I15-474}\quad \textbf{JEL:} \texttt{I15}\par
% BEGIN ECONOMICS STATEMENT
\label{problem:OP-0037}
\label{atlas:prob:D7}
\noindent\textbf{Original atlas ID:} D7 (entry 37). \textbf{Classification:} Editorial health-policy, demographic, and long-run equilibrium problem.

\noindent\textbf{Original atlas question:} How large are the long-run effects of early-life disease and nutrition on education, fertility, productivity and growth?

\subsection*{Definitions and assumptions}
Fix a baseline birth or child cohort, health-policy versions $\pi$, and horizon $H$. Policy specifies disease prevention, nutrition, eligibility, uptake support, and timing; different interventions need not affect the same biological pathway. Let $A_{it}(\pi)$ indicate survival, $E_{it}(\pi)$ completed schooling years through date $t$, $F_{it}(\pi)$ a defined fertility outcome, and $Y_{it}(\pi)$ real earnings. To define a baseline-cohort earnings target, set earnings to zero after death under an explicit accounting convention and report survival separately. Cumulative schooling and fertility remain defined after death by retaining their attained predeath values. This earnings convention does not assign a welfare value to death. Permit disease-transmission and household spillovers by indexing outcomes by full allocation vectors or a declared exposure mapping.
\subsection*{Formal research question}
For each policy pair and date $t\in\{0,\ldots,H\}$, define
\begin{align*}
 \theta^A_t&=\mathbb E[A_{it}(\pi)-A_{it}(\pi_0)],\\
 \theta^E_t&=\mathbb E[E_{it}(\pi)-E_{it}(\pi_0)],\\
 \theta^F_t&=\mathbb E[F_{it}(\pi)-F_{it}(\pi_0)],\\
 \theta^Y_H&=\mathbb E\sum_{t=0}^H\beta^t
                    [Y_{it}(\pi)-Y_{it}(\pi_0)].
\end{align*}
Identify or sharply bound the fully indexed vector
\[
 \Theta_H(\pi,\pi_0)
    =\bigl((\theta^A_t,\theta^E_t,\theta^F_t)_{t=0}^H,\theta^Y_H\bigr),
\]
with $\beta\in(0,1]$ and the stated outcome conventions at death. Determine transport across disease environments, nutritional deficits, and implementation systems. Separately, in a model of mortality, fertility, migration, schooling, and capital accumulation, characterize national output and welfare responses; specify whether output is total, per baseline person, or per contemporaneous resident.
\subsection*{Interpretation and scope}
A health intervention's total effect is different from the causal effect of a one-unit improvement in a health index. An uptake instrument needs exclusions that may fail when access changes information, attendance, or household resources. Conditioning on adult survivors can compare different selected populations; baseline-cohort outcomes or explicitly bounded survivor effects are needed. Growth rates and output per capita can change through both the numerator and demographic denominator. Randomization can identify long-run policy effects when baseline assignment is maintained, but requires tracking, treatment consistency, and an interference model; the passage of time alone does not invalidate randomization. Disease-control evidence cannot establish every nutritional or macroeconomic feedback without additional assumptions.
\subsection*{Source audit and status}
[S1] concerns deworming externalities and [S3] long-run outcomes. ``Bleakley (2010)'' is ambiguous between the verified review [S2] and malaria article [S4]. Added [S5] verifies twenty-year deworming follow-up, so adult economic effects are not wholly unmeasured. The cross-intervention, demographic, and national-equilibrium extension is editorial; no universal unresolved 2026 status is asserted.

\subsection*{References}
\begin{enumerate}[label=\textnormal{[S\arabic*]},leftmargin=*]
\item Edward Miguel and Michael Kremer (2004). \emph{Worms: Identifying Impacts on Education and Health in the Presence of Treatment Externalities}. Econometrica 72(1), 159--217. \url{https://doi.org/10.1111/j.1468-0262.2004.00481.x}. \textit{Locator:} Primary publisher, working-paper, or author record: bibliographic information and abstract; full-text theorem verification is not claimed. \textit{Role:} School-based deworming and treatment-externality evidence.
\item Hoyt Bleakley (2010). \emph{Health, Human Capital, and Development}. Annual Review of Economics 2, 283--310. \url{https://doi.org/10.1146/annurev.economics.102308.124436}. \textit{Locator:} Primary publisher, working-paper, or author record: bibliographic information and abstract; full-text theorem verification is not claimed. \textit{Role:} Verified broad review candidate for the health-development shorthand. Author-year is ambiguous: a second 2010 article on malaria eradication is also a plausible intended source. Both are supplied rather than asserting one exact original identity.
\item Sarah Baird, Joan Hamory Hicks, Michael Kremer, and Edward Miguel (2016). \emph{Worms at Work: Long-Run Impacts of a Child Health Investment}. The Quarterly Journal of Economics 131(4), 1637--1680. \url{https://academic.oup.com/qje/article-abstract/131/4/1637/2468871}. \textit{Locator:} Primary publisher, working-paper, or author record: bibliographic information and abstract; full-text theorem verification is not claimed. \textit{Role:} Long-run follow-up of a particular childhood deworming intervention.
\item Hoyt Bleakley (2010). \emph{Malaria Eradication in the Americas: A Retrospective Analysis of Childhood Exposure}. American Economic Journal: Applied Economics 2(2), 1--45. \url{https://doi.org/10.1257/app.2.2.1}. \textit{Locator:} Primary publisher, working-paper, or author record: bibliographic information and abstract; full-text theorem verification is not claimed. \textit{Role:} Second verified candidate; historical disease-control evidence.
\item Joan Hamory, Edward Miguel, Michael Walker, Michael Kremer, and Sarah Baird (2020). \emph{Twenty Year Economic Impacts of Deworming}. NBER Working Paper 27611; publisher record reports a 2021 journal version. \url{https://www.nber.org/papers/w27611}. \textit{Locator:} Primary publisher, working-paper, or author record: bibliographic information and abstract; full-text theorem verification is not claimed. \textit{Role:} Extends follow-up to twenty years, so the atlas must not imply that all adult effects remain unmeasured.
\end{enumerate}

\subsection*{Common conventions: Source mathematical conventions}

Unless an entry states otherwise, agent and firm sets are finite. State and action spaces are measurable, strategies and outcome maps are measurable, probability laws are specified, and expectations exist for the targets under discussion. Infinite-horizon discount factors lie in $(0,1)$ and discounted objectives are finite. Norms, tolerances and complexity input sizes must be specified for computational comparisons. Constants or primitives that appear locally are inputs, not empirical facts asserted by this catalogue.

\textbf{Identification.} A statistical model $m\in\mathcal M$ implies an observable law $P_m$ and target $\theta(m)$. At law $P$, its identified set is
\[
\Theta(P;\mathcal M)=\{\theta(m):m\in\mathcal M,\ P_m=P\}.
\]
Point identification holds when this set is a singleton, or equivalently when $P_{m_1}=P_{m_2}$ implies $\theta(m_1)=\theta(m_2)$ for all compatible models. Sharp bounds mean bounds attained, or approached when the boundary is not attained, by compatible models. Writing this set is a definition, not a solution: the substantive task is to establish informative restrictions and derive or estimate the set. An unrestricted-confounding model may give only trivial bounds.

\textbf{Inference.} Identification, estimability and valid uncertainty quantification are separate questions. Fix $\alpha\in(0,1)$. A confidence procedure $C_n$ over a declared law class $\mathcal P$ has asymptotic uniform coverage at level $1-\alpha$ when
\[
\liminf_{n\to\infty}\inf_{P\in\mathcal P}\Pr_P\{\theta(P)\in C_n\}\geq1-\alpha.
\]
For set-identified targets the coverage event must be defined separately, for example coverage of the full identified set. If an entry uses identification-indexed models instead of uniquely indexed laws, the same coverage convention is applied over those models. Pointwise limits do not establish uniform validity, and asymptotic validity does not guarantee finite-sample precision. Sample sizes, dependence structures and weak-identification sequences are part of each application.

\textbf{Counterfactuals.} $Y_i(a)$ is a potential outcome under a specified intervention, including its timing, implementation and versions. It is not $Y_i$ conditional on voluntarily choosing $a$. A full intervention vector permits interference; shorthand scalar treatment notation does not silently impose no interference. Assignment, selection, attrition, anticipation and measurement must be specified. A claim about an instrument's local effect is not automatically a population policy effect.

\textbf{Economic equilibrium.} A counterfactual model includes best responses, constraints, feasibility, market clearing, state transitions and expectations consistent with the stated information. When several equilibria exist, either a declared selector is an input or the target is set-valued. Comparisons across policy regimes must use the stated selection convention. Reduced-form relationships are not presumed invariant after an equilibrium-changing intervention.

\textbf{Optimization and welfare.} Optimization is over the stated feasible set. If attainment is not established, $\sup$ or $\inf$ and $\varepsilon$-optimal choices replace an asserted maximizer or minimizer. For example, with value $V=\sup_{a\in\mathcal A}W(a)<\infty$, an $\varepsilon$-optimal set is $\{a\in\mathcal A:W(a)\geq V-\varepsilon\}$ for $\varepsilon>0$. Benefits, costs, risk and health or environmental outcomes can be aggregated only using declared units and conversion weights. Interpersonal utility weights, the treatment of fiscal transfers and foreign profits, discounting, and the behavioral welfare criterion are explicit inputs. A comparative-static derivative additionally requires existence and appropriate differentiability of the selected counterfactual.

\textbf{Sources.} Local labels [S1], [S2], and so forth restart in every question. Locators identify the relevant abstract, model, section or result at the level actually checked; a bibliographic or abstract-level verification is not represented as inspection of an entire proof. Working papers are distinguished from later publications and changing versions. Source summaries are paraphrased. All URL checks and editorial formulations refer to the audit date stated above.
% END ECONOMICS STATEMENT
\end{document}
